% !TEX root = ../main.tex
\begin{ouverture}[transferts]{20}{Transferts thermiques et bilans}{Partie 3 · L'énergie : conversions et transferts}
{Conduction, convection, rayonnement\enspace·\enspace Résistance thermique\enspace·\enspace Bilan radiatif de la Terre\enspace·\enspace Loi de Newton}
\objectif{Caractériser qualitativement les trois modes de transfert thermique : conduction, convection, rayonnement.}
\objectif{Exploiter la relation entre flux thermique, résistance thermique et écart de température.}
\objectif{Effectuer un bilan d'énergie pour estimer la température moyenne de la Terre, la loi de Stefan-Boltzmann étant
donnée ; discuter de l'influence de l'albédo et de l'effet de serre.}
\objectif{Établir et résoudre l'équation différentielle qui régit la température d'un système au contact d'un thermostat.}
\objectif{Modéliser l'évolution de la température par la méthode d'Euler en Python.}
\prerequis{Premier principe, capacité thermique (chapitre~19).}
\prerequis{Bilan radiatif terrestre, albédo (enseignement scientifique de première).}
\prerequis{Équations différentielles $y'=a\,y+b$ (mathématiques) ; méthode d'Euler (chapitres~14 et~17).}
\end{ouverture}
\essentiel
\essSection{1}{Les modes de transfert thermique}
\begin{definition}[Conduction, convection, rayonnement]
\begin{itemize}
\item \textbf{Conduction} : l'agitation thermique se transmet de proche en proche par chocs entre entités voisines,
sans déplacement de matière. C'est le mode principal dans les solides.
\item \textbf{Convection} : un fluide chauffé se dilate, devient moins dense et s'élève ; il est remplacé par du fluide
plus froid. Le transfert accompagne un déplacement de matière.
\item \textbf{Rayonnement} : tout corps émet des ondes électromagnétiques, d'autant plus intenses qu'il est chaud, et
absorbe une partie de celles qu'il reçoit. Ce transfert se produit même dans le vide.
\end{itemize}
\end{definition}
\schema{20-modes}
\essSection{2}{Flux thermique et résistance thermique}
\begin{loi}[Flux thermique et résistance thermique]
Le \textbf{flux thermique} $\Phi=\dfrac{Q}{\Delta t}$ à travers une paroi est l'énergie transférée par unité de temps, en
watts. En régime permanent, le flux qui traverse une paroi dont les faces sont aux températures $T_1>T_2$ est
proportionnel à l'écart de température :
\[\Phi=\frac{T_1-T_2}{R_\mathrm{th}},\]
où $R_\mathrm{th}$ est la résistance thermique de la paroi, en \unit{K.W^{-1}}. Pour une paroi plane d'épaisseur $e$, de surface
$S$, faite d'un matériau de conductivité thermique $\lambda$ (en \unit{W.m^{-1}.K^{-1}}), elle s'exprime par
$R_\mathrm{th}=\dfrac{e}{\lambda\,S}$. Lorsque plusieurs couches sont traversées successivement par le même flux, leurs
résistances thermiques s'ajoutent.
\end{loi}
\begin{exemple}
Un mur de béton ($\lambda=\qty{1,75}{W.m^{-1}.K^{-1}}$) de \qty{20}{cm} d'épaisseur et de \qty{10}{m^2} a une résistance thermique
$R=\dfrac{0{,}20}{1{,}75\times10}=\qty{0,0114}{K.W^{-1}}$ : pour \qty{20}{\degreeCelsius} d'écart, il laisse passer
\qty{1,75e3}{W}. Recouvert de \qty{10}{cm} de laine de verre ($\lambda=\qty{0,040}{W.m^{-1}.K^{-1}}$, $R=\qty{0,25}{K.W^{-1}}$), il ne
laisse plus passer que \qty{77}{W}, 23 fois moins.
\end{exemple}
\schema{20-mur}
\essSection{3}{Le bilan radiatif de la Terre}
\begin{loi}[Loi de Stefan-Boltzmann]
Un corps de température $T$ qui absorbe tout rayonnement émet une puissance par unité de surface
$\dfrac{P}{S}=\sigma\,T^4$, avec $\sigma=\qty{5,67e-8}{W.m^{-2}.K^{-4}}$. Cette loi est fournie dans les exercices.
\end{loi}
\begin{savoirfaire}[Estimer la température de la Terre sans atmosphère]
Le Soleil apporte à la Terre la puissance $S_0=\qty{1361}{W.m^{-2}}$ par mètre carré perpendiculaire à ses rayons. La Terre
intercepte ce rayonnement sur un disque de surface $\pi R^2$ et en réfléchit la fraction $A=0{,}30$ (l'\textbf{albédo}) ;
elle émet, en infrarouge, sur toute sa surface $4\pi R^2$. À l'équilibre, puissance absorbée et puissance émise sont
égales :
\[(1-A)\,S_0\,\pi R^2=\sigma\,T^4\times4\pi R^2\qquad\text{d'où}\qquad T=\left(\frac{(1-A)\,S_0}{4\,\sigma}\right)^{1/4}=\qty{255}{K},\]
soit environ \qty{-18}{\degreeCelsius}, alors que la température moyenne mesurée au sol vaut \qty{15}{\degreeCelsius}.
\end{savoirfaire}
\schema{20-terre}
\begin{propriete}[Albédo et effet de serre]
Une surface claire (neige, glace, nuages) réfléchit davantage le rayonnement solaire : un albédo plus élevé abaisse
la température d'équilibre. L'atmosphère, presque transparente au rayonnement solaire, absorbe en partie le
rayonnement infrarouge émis par le sol, grâce à la vapeur d'eau, au dioxyde de carbone et au méthane, et en réémet
une partie vers le sol : c'est l'\textbf{effet de serre}, qui porte la température moyenne de \qty{-18}{\degreeCelsius} à
\qty{15}{\degreeCelsius}. Son renforcement par les émissions humaines réchauffe le climat.
\end{propriete}
\essSection{4}{Évolution de la température d'un système}
\begin{loi}[Loi phénoménologique de Newton]
Un système de température $T$, de surface $S$, au contact d'un fluide extérieur de température $T_\mathrm{ext}$ constante
(un thermostat), reçoit le flux thermique
\[\Phi=h\,S\,(T_\mathrm{ext}-T),\]
où $h$ est le coefficient d'échange (en \unit{W.m^{-2}.K^{-1}}). Cette loi est fournie dans les exercices.
\end{loi}
\begin{savoirfaire}[Établir et résoudre l'équation différentielle]
Pour un système incompressible de capacité thermique $C$ qui n'échange pas de travail, le premier principe appliqué
entre $t$ et $t+\mathrm{d}t$ s'écrit $C\,\mathrm{d}T=\Phi\,\mathrm{d}t$, soit
\[\frac{\mathrm{d}T}{\mathrm{d}t}=-\frac{1}{\tau}\,(T-T_\mathrm{ext})\qquad\text{avec}\qquad\tau=\frac{C}{h\,S}.\]
Cette équation linéaire du premier ordre, à coefficients constants et second membre constant, a pour solution, avec
$T(0)=T_0$ :
\[T(t)=T_\mathrm{ext}+(T_0-T_\mathrm{ext})\,\mathrm{e}^{-t/\tau}.\]
Le \textbf{temps caractéristique} $\tau$ mesure la rapidité de l'évolution : à $t=\tau$, il reste 37~\% de l'écart initial ;
à $t=5\tau$, moins de 1~\%.
\end{savoirfaire}
\schema{20-newton}
\begin{codecote}[0.6]
\begin{pythoncol}
from math import exp
Text, T0, tau = 20.0, 80.0, 1500.0  # °C, s
T, t, dt = T0, 0.0, 10.0
while t < 3600:
    T = T + (Text - T) / tau * dt
    t = t + dt
print(round(T, 2),
      round(Text + (T0-Text)*exp(-t/tau), 2))
\end{pythoncol}
\begin{consolecol}
25.4 25.44
\end{consolecol}
\cote
\fontsize{9.5}{12.5}\selectfont
\capacite{Capacité numérique}\par\smallskip
Le programme applique la méthode d'Euler à l'équation $\dfrac{\mathrm{d}T}{\mathrm{d}t}=\dfrac{T_\mathrm{ext}-T}{\tau}$ : à chaque pas
$\mathrm{d}t$, la température varie de $\dfrac{T_\mathrm{ext}-T}{\tau}\,\mathrm{d}t$. Au bout d'une heure, une boisson à
\qty{80}{\degreeCelsius} atteint \qty{25,4}{\degreeCelsius}, contre \qty{25,44}{\degreeCelsius} avec la solution exacte.
\end{codecote}
\begin{attention}
Ne pas confondre le flux thermique $\Phi$, en watts, et le transfert thermique $Q=\Phi\,\Delta t$, en joules. Un écart de
température a la même valeur en kelvins et en degrés Celsius, mais la loi de Stefan-Boltzmann exige la température
en kelvins.
\end{attention}
\begin{experience}[Suivre et modéliser un refroidissement]
On remplit un récipient d'eau chaude et l'on enregistre sa température toutes les trente secondes avec une sonde reliée
à un microcontrôleur, la pièce étant à température constante. On trace $T$ en fonction de $t$, puis
$\ln(T-T_\mathrm{ext})$ en fonction de $t$ : on obtient une droite de coefficient directeur $-1/\tau$, ce qui valide le
modèle de Newton et donne le temps caractéristique.
\end{experience}
\begin{rediger}
Pour établir l'équation différentielle, définir le système, écrire le premier principe entre $t$ et $t+\mathrm{d}t$,
$\mathrm{d}U=\delta Q$, puis remplacer $\mathrm{d}U=C\,\mathrm{d}T$ et $\delta Q=\Phi\,\mathrm{d}t$. Identifier $\tau$, vérifier que la
solution proposée convient et utiliser la condition initiale pour fixer la constante.
\end{rediger}
\exercices

\rubrique{Découverte}
\begin{blocLivre}{ex-20-1}
\begin{exercice}[1]{Trois modes de transfert}
\begin{questions}
\item Une cuillère métallique plongée dans un thé chaud devient brûlante à son extrémité. Quel mode de transfert est en
jeu ?
\item Un radiateur placé sous une fenêtre chauffe toute la pièce. Expliquer, en précisant le mode de transfert.
\item Comment l'énergie du Soleil nous parvient-elle à travers l'espace ?
\item Une bouteille isotherme a une double paroi de verre argentée, entre lesquelles on a fait le vide. Quel mode de
transfert chaque élément de sa conception limite-t-il ?
\end{questions}
\end{exercice}
\end{blocLivre}
\begin{blocLivre}{ex-20-2}
\begin{exercice}[1]{La porte d'entrée}
Une porte en bois de \qty{2,0}{m^2} et de \qty{4,0}{cm} d'épaisseur sépare une maison à \qty{20}{\degreeCelsius} de l'extérieur à
\qty{1}{\degreeCelsius}. On admet que ses faces sont aux températures de l'intérieur et de l'extérieur.
\begin{donnees}
$\lambda_\mathrm{bois}=\qty{0,15}{W.m^{-1}.K^{-1}}$ ; $R_\mathrm{th}=\dfrac{e}{\lambda\,S}$.
\end{donnees}
\begin{questions}
\item Calculer la résistance thermique de la porte.
\item Calculer le flux thermique qui la traverse. Dans quel sens ?
\item Quelle énergie la maison perd-elle par cette porte en 24~heures ? L'exprimer en kilowattheures.
\item Quelle épaisseur de bois diviserait ce flux par deux ?
\end{questions}
\end{exercice}
\end{blocLivre}
\begin{blocLivre}{ex-20-3}
\begin{exercice}[1]{Rayonner selon sa température}
\begin{donnees}
Loi de Stefan-Boltzmann : $\dfrac{P}{S}=\sigma\,T^4$, avec $\sigma=\qty{5,67e-8}{W.m^{-2}.K^{-4}}$.
\end{donnees}
\begin{questions}
\item Calculer la puissance émise par mètre carré par la surface du Soleil, à \qty{5772}{K}.
\item Même question pour le sol terrestre à \qty{288}{K}, puis pour la peau humaine à \qty{306}{K}.
\item Combien de fois la surface du Soleil émet-elle plus que celle de la Terre, par mètre carré ?
\item Par quel facteur la puissance rayonnée est-elle multipliée si la température thermodynamique double ?
\end{questions}
\end{exercice}
\end{blocLivre}
\begin{blocLivre}{ex-20-4}
\begin{exercice}[1]{Le thé refroidit}
Une tasse de thé, posée dans une pièce à \qty{20}{\degreeCelsius}, refroidit selon la loi de Newton :
$T(t)=T_\mathrm{ext}+(T_0-T_\mathrm{ext})\,\mathrm{e}^{-t/\tau}$.
\schema{20-ex04}
\begin{questions}
\item Lire la température initiale du thé. Vers quelle valeur la température tend-elle ?
\item Déterminer graphiquement le temps caractéristique $\tau$ par deux méthodes : la tangente à l'origine, puis la
valeur de $T$ pour laquelle il reste 37~\% de l'écart initial.
\item Calculer la température du thé au bout de $2\tau$, puis de $3\tau$. Vérifier sur la courbe.
\item On aime boire le thé à \qty{60}{\degreeCelsius}. Combien de temps faut-il attendre ?
\end{questions}
\end{exercice}
\end{blocLivre}
\begin{blocLivre}{ex-20-5}
\begin{exercice}[1]{L'albédo des surfaces}
Le rayonnement solaire apporte en moyenne \qty{340}{W.m^{-2}} au sommet de l'atmosphère. On compare des surfaces dont
l'albédo vaut : neige fraîche 0,85 ; forêt 0,15 ; océan 0,06.
\begin{questions}
\item Définir l'albédo d'une surface.
\item Calculer la puissance absorbée par mètre carré par chaque surface, en négligeant l'atmosphère.
\item Quelle surface s'échauffe le plus au soleil ? le moins ?
\item La fonte de la banquise remplace de la glace par de l'océan. Expliquer pourquoi ce phénomène accentue le
réchauffement des régions polaires.
\end{questions}
\end{exercice}
\end{blocLivre}
\rubrique{Entraînement}
\begin{blocLivre}{ex-20-6}
\begin{exercice}[2]{Simple ou double vitrage}
Pour une fenêtre de \qty{1,5}{m^2}, les résistances thermiques totales, couches d'air au contact des vitres comprises,
valent : simple vitrage \qty{0,115}{K.W^{-1}} ; double vitrage \qty{0,238}{K.W^{-1}} ; double vitrage performant (gaz argon,
couche peu émissive) \qty{0,606}{K.W^{-1}}.
\begin{questions}
\item Calculer le flux thermique perdu par chaque fenêtre pour un écart de \qty{20}{\degreeCelsius} entre intérieur et
extérieur.
\item Les fabricants indiquent un coefficient $U=\dfrac{1}{R_\mathrm{th}\,S}$, en \unit{W.m^{-2}.K^{-1}}. Calculer $U$ pour le simple
vitrage et interpréter.
\item Pendant l'hiver, l'écart moyen de température vaut \qty{12}{\degreeCelsius} durant \num{3600} heures. Quelle énergie
économise-t-on en remplaçant le simple vitrage par le double vitrage performant ?
\item Quelle économie cela représente-t-il à 0,20~€ le kilowattheure ?
\end{questions}
\end{exercice}
\end{blocLivre}
\begin{blocLivre}{ex-20-7}
\begin{exercice}[2]{Isoler les combles}
Un plafond de \qty{80}{m^2} sépare une maison chauffée à \qty{20}{\degreeCelsius} de combles non chauffés. Sans isolant, sa
résistance thermique totale vaut \qty{3,0e-3}{K.W^{-1}}. On le recouvre de \qty{30}{cm} de laine de verre.
\begin{donnees}
$\lambda_\mathrm{laine}=\qty{0,040}{W.m^{-1}.K^{-1}}$ ; $R_\mathrm{th}=\dfrac{e}{\lambda\,S}$ ; les résistances de couches traversées
successivement s'ajoutent.
\end{donnees}
\begin{questions}
\item Calculer la résistance thermique de la couche de laine de verre.
\item Pour un écart de \qty{20}{\degreeCelsius}, calculer le flux perdu avant et après isolation.
\item Les aides publiques à la rénovation exigent une résistance thermique surfacique $R_\mathrm{th}\times S=e/\lambda$ d'au
moins \qty{7}{m^2.K.W^{-1}}. L'isolation est-elle éligible ? Quelle épaisseur minimale faut-il ?
\item Avec un écart moyen de \qty{12}{\degreeCelsius} pendant \num{3600} heures, quelle énergie l'isolation fait-elle
économiser ?
\end{questions}
\end{exercice}
\end{blocLivre}
\begin{blocLivre}{ex-20-8}
\begin{exercice}[2]{Modéliser un refroidissement en Python}
On reprend le programme de l'essentiel.
\begin{questions}
\item \capacite{Capacité numérique}Quelle relation la ligne \code{T = T + (Text - T) / tau * dt} traduit-elle ? D'où
vient-elle ?
\item Pourquoi la valeur obtenue par la méthode d'Euler diffère-t-elle légèrement de la valeur exacte ? Comment réduire
cet écart ?
\item Modifier le programme pour qu'il affiche l'instant où la température passe sous \qty{30}{\degreeCelsius}. Retrouver ce
résultat par le calcul.
\item Écrire les lignes qui tracent la température en fonction du temps, en minutes.
\end{questions}
\end{exercice}
\end{blocLivre}
\begin{blocLivre}{ex-20-9}
\begin{exercice}[2]{La température de la Terre sans atmosphère}
\begin{donnees}
Constante solaire $S_0=\qty{1361}{W.m^{-2}}$ ; albédo terrestre $A=0{,}30$ ; $\sigma=\qty{5,67e-8}{W.m^{-2}.K^{-4}}$ ; loi de
Stefan-Boltzmann : $P/S=\sigma\,T^4$.
\end{donnees}
\begin{questions}
\item Exprimer la puissance solaire absorbée par la Terre, de rayon $R$, et la puissance qu'elle rayonne.
\item En déduire sa température d'équilibre, en kelvins puis en degrés Celsius.
\item Calculer la température d'équilibre d'une Terre entièrement couverte de glace ($A=0{,}60$), puis d'une Terre
parfaitement noire ($A=0$).
\item La température moyenne réelle vaut \qty{15}{\degreeCelsius}. Quel phénomène, absent de ce modèle, explique l'écart ?
\end{questions}
\end{exercice}
\end{blocLivre}
\begin{blocLivre}{ex-20-10}
\begin{exercice}[2]{Froid près d'une fenêtre}
Une personne de surface corporelle \qty{1,8}{m^2} a une peau à \qty{33}{\degreeCelsius}. Les murs de la pièce sont à
\qty{20}{\degreeCelsius}. On assimile la peau et les murs à des corps qui absorbent tout rayonnement.
\begin{donnees}
$\sigma=\qty{5,67e-8}{W.m^{-2}.K^{-4}}$ ; loi de Stefan-Boltzmann : $P/S=\sigma\,T^4$.
\end{donnees}
\begin{questions}
\item Calculer la puissance rayonnée par la peau.
\item La peau absorbe le rayonnement des murs, soit $\sigma\,S\,T_\mathrm{mur}^4$. Calculer la puissance nette perdue par
rayonnement.
\item Même question face à une grande baie vitrée dont la surface intérieure est à \qty{10}{\degreeCelsius}.
\item Au repos, le corps produit environ \qty{100}{W}. Expliquer la sensation de froid près d'une baie vitrée, même si
l'air de la pièce est à \qty{20}{\degreeCelsius}.
\end{questions}
\end{exercice}
\end{blocLivre}
\begin{blocLivre}{ex-20-11}
\begin{exercice}[2]{Rafraîchir une canette}
Une canette contient \qty{330}{g} de soda, assimilé à de l'eau, à \qty{25}{\degreeCelsius}. L'aluminium de la canette a une capacité
thermique de \qty{13}{J.K^{-1}}. On la place dans un réfrigérateur à \qty{4}{\degreeCelsius}.
\begin{donnees}
$c_\mathrm{eau}=\qty{4180}{J.kg^{-1}.K^{-1}}$ ; surface de la canette : $S=\qty{0,030}{m^2}$ ; coefficient d'échange avec l'air :
$h=\qty{10}{W.m^{-2}.K^{-1}}$ ; loi de Newton : $\Phi=h\,S\,(T_\mathrm{ext}-T)$.
\end{donnees}
\begin{questions}
\item Établir l'équation différentielle vérifiée par la température $T$ de la canette.
\item Calculer la capacité thermique de la canette pleine, puis le temps caractéristique $\tau$.
\item Au bout de combien de temps le soda atteint-il \qty{8}{\degreeCelsius} ?
\item Même question dans un congélateur à \qty{-18}{\degreeCelsius}, le coefficient d'échange étant le même. Pourquoi ne faut-il pas
y oublier la canette ?
\end{questions}
\end{exercice}
\end{blocLivre}
\rubrique{Objectif bac}
\begin{blocLivre}{ex-20-12}
\begin{exercice}[3]{Un modèle de l'effet de serre}
\begin{docu}{1}{Une atmosphère en une couche}
On modélise l'atmosphère par une couche transparente au rayonnement solaire, qui absorbe la fraction $\varepsilon$ du
rayonnement infrarouge émis par le sol et en réémet la moitié vers le sol, l'autre moitié vers l'espace. Le bilan
d'énergie du sol conduit alors à $\sigma\,T_\mathrm{s}^4\left(1-\dfrac{\varepsilon}{2}\right)=\sigma\,T_\mathrm{e}^4$, où $T_\mathrm{e}$ est la
température d'équilibre sans atmosphère.
\end{docu}
\begin{docu}{2}{Valeurs}
Sans atmosphère, $T_\mathrm{e}=\qty{255}{K}$ (valeur arrondie de $\left(\dfrac{(1-A)\,S_0}{4\sigma}\right)^{1/4}=\qty{254,6}{K}$). Température
moyenne mesurée au sol : \qty{288}{K}.
\end{docu}
\begin{questions}
\item Que vaut $T_\mathrm{s}$ si $\varepsilon=0$ ? Interpréter.
\item Calculer $T_\mathrm{s}$ pour une atmosphère qui absorberait tout le rayonnement infrarouge ($\varepsilon=1$).
\item \initiative Quelle valeur de $\varepsilon$ rend compte de la température moyenne mesurée ? On utilisera
$T_\mathrm{e}=\qty{254,6}{K}$.
\item L'augmentation de la concentration en dioxyde de carbone fait croître $\varepsilon$. Calculer la hausse de température
si $\varepsilon$ passe de 0,78 à 0,80, et conclure.
\end{questions}
\end{exercice}
\end{blocLivre}
\begin{blocLivre}{ex-20-13}
\begin{exercice}[3]{Mars et Vénus}
\begin{docu}{1}{Deux voisines de la Terre}
\begin{center}
\begin{tabular}{lccc}
\toprule
Planète & distance au Soleil (ua) & albédo & température moyenne mesurée \\
Mars & 1,524 & 0,25 & \qty{210}{K} \\
Vénus & 0,723 & 0,76 & \qty{737}{K} \\
\bottomrule
\end{tabular}
\end{center}
\end{docu}
\begin{docu}{2}{Leurs atmosphères}
L'atmosphère de Mars, très ténue (pression au sol : \qty{6}{hPa}), contient surtout du dioxyde de carbone. Celle de Vénus,
presque entièrement composée de dioxyde de carbone, exerce au sol une pression de \qty{92}{bar}.
\end{docu}
\begin{donnees}
Constante solaire à \qty{1}{ua} : \qty{1361}{W.m^{-2}} ; la puissance reçue varie comme l'inverse du carré de la distance ;
$\sigma=\qty{5,67e-8}{W.m^{-2}.K^{-4}}$.
\end{donnees}
\begin{questions}
\item Calculer la puissance solaire reçue par mètre carré au niveau de Mars et au niveau de Vénus.
\item Établir l'expression de la température d'équilibre d'une planète sans effet de serre, puis la calculer pour Mars et
pour Vénus.
\item \initiative Comparer aux températures mesurées et interpréter à l'aide du document~2.
\item Vénus est plus éloignée du Soleil que Mercure, mais beaucoup plus chaude. Expliquer.
\end{questions}
\end{exercice}
\end{blocLivre}
\begin{blocLivre}{ex-20-14}
\begin{exercice}[3]{La bouteille isotherme}
\begin{docu}{1}{Un café qui reste chaud}
On verse \qty{0,50}{L} de café à \qty{90}{\degreeCelsius} dans une bouteille isotherme, dans une pièce à \qty{20}{\degreeCelsius}. Six
heures plus tard, il est encore à \qty{70}{\degreeCelsius}. Dans une tasse ordinaire, le temps caractéristique de refroidissement
vaut environ \qty{25}{min}.
\end{docu}
\begin{docu}{2}{Modèle}
On admet que la température suit la loi $T(t)=T_\mathrm{ext}+(T_0-T_\mathrm{ext})\,\mathrm{e}^{-t/\tau}$, avec $\tau=\dfrac{C}{h\,S}$. La
capacité thermique du café et de la paroi intérieure vaut $C=\qty{2,3e3}{J.K^{-1}}$.
\end{docu}
\begin{questions}
\item Calculer le temps caractéristique de la bouteille isotherme.
\item Combien de fois la bouteille ralentit-elle le refroidissement par rapport à une tasse ?
\item \initiative Au bout de combien de temps le café passe-t-il sous \qty{60}{\degreeCelsius} ? Quelle est sa température au bout de
12~heures ?
\item Calculer le produit $h\,S$ pour la bouteille et expliquer, à l'aide des modes de transfert, pourquoi il est si
faible.
\end{questions}
\end{exercice}
\end{blocLivre}
\begin{blocLivre}{ex-20-15}
\begin{exercice}[3]{L'heure du crime}
\begin{docu}{1}{Les constatations}
Un corps est découvert dans une pièce dont la température est restée égale à \qty{18}{\degreeCelsius}. À 23~h, le médecin
légiste mesure une température corporelle de \qty{30}{\degreeCelsius} ; à 1~h du matin, de \qty{26}{\degreeCelsius}.
\end{docu}
\begin{docu}{2}{Le modèle}
On admet que la température du corps, égale à \qty{37}{\degreeCelsius} au moment du décès, suit ensuite la loi de refroidissement
de Newton, avec un temps caractéristique constant. En pratique, les médecins légistes utilisent des modèles plus
élaborés.
\end{docu}
\begin{questions}
\item Écrire l'expression de $T(t)$ et montrer que $\dfrac{T(t_1)-T_\mathrm{ext}}{T(t_2)-T_\mathrm{ext}}=\mathrm{e}^{(t_2-t_1)/\tau}$.
\item En déduire le temps caractéristique de refroidissement du corps.
\item \initiative Estimer l'heure du décès.
\item Pourquoi les médecins légistes considèrent-ils ce résultat comme une simple estimation ?
\end{questions}
\end{exercice}
\end{blocLivre}
